\NSPage{145}%
\begin{EESegment}{EE-TGT-P145-001}
The function \NSi{NS-F-P145-001}{\Pi_0} is holomorphic. For
\NSi{NS-F-P145-002}{\Lambda\geq 1}, define
\end{EESegment}
\NSFormula{NS-F-P145-003}{display}{B.3}
\begin{equation}
  \zeta_*=-\frac{LH_*}{H_*^2+\sigma_*^2},
  \qquad \widetilde\xi_0=\Lambda\zeta_*,
  \qquad \phi_*=\exp\!\left(\Lambda\int_0^\eta \zeta_*(w)\,dw\right).
  \tag{B.3}\label{eq:B.3}
\end{equation}
\begin{EESegment}{EE-TGT-P145-002}
The integral is well defined throughout \NSi{NS-F-P145-004}{\Omega}. On the real interval,
\NSi{NS-F-P145-005}{\phi_*>0}. The amplitude \NSi{NS-F-P145-006}{\mathsf C} is chosen after
\NSi{NS-F-P145-007}{\Lambda}.
\end{EESegment}

\begin{EESegment}{EE-TGT-P145-003}
\subsection{A space for analytic coefficients.}\label{sec:B.2}
The profile equations \eqref{eq:4.13} use parameter derivatives as well as radial derivatives. Their highest radial
derivatives are multiplied by the radial variable. At \NSi{NS-F-P145-008}{X=0} that multiplier
vanishes, so the equations cannot be solved there for their highest radial derivatives in the
usual way. The vanishing is only linear, and the smooth power-series coefficients can still be
solved one at a time by the radial inverse in \eqref{eq:B.5}. That is what regular singular point
means here. This subsection sets up a weighted
coefficient space \NSi{NS-F-P145-009}{\mathcal B_\rho} for the integral form of those equations. Its
norm is given in \eqref{eq:B.4}. In this space, increasing the radial degree makes up for taking a
parameter derivative. Lemma~\ref{lem:B.1} gives the bounds needed for multiplication and radial
inversion, the step that solves for a function in the radial variable. Those bounds also cover the
mixed nonlinear terms used in the contraction argument for Proposition~\ref{prop:B.2}.
\end{EESegment}

\begin{EESegment}{EE-TGT-P145-004}
For the analytic construction at the axis, use the scalar radial variable
\NSi{NS-F-P145-010}{Y=\Lambda X}. Fix a small radius \NSi{NS-F-P145-011}{\rho>0}. If
\NSi{NS-F-P145-012}{F(Y,\eta)=\sum_{\alpha\geq0}F_\alpha(\eta)Y^\alpha}, define
\end{EESegment}
\NSFormula{NS-F-P145-013}{display}{B.4}
\begin{equation}
  \mathfrak a_{\alpha\beta}
  =\frac{20^{-\alpha}\rho^{-\beta}\beta!\binom{\alpha+\beta}{\beta}}
         {(\alpha+1)^2(\beta+1)^2},
  \qquad
  \lVert F\rVert_\rho
  =\sup_{\alpha,\beta\geq0}\sup_{\eta\in I}
    \frac{|\partial_\eta^\beta F_\alpha(\eta)|}{\mathfrak a_{\alpha\beta}}.
  \tag{B.4}\label{eq:B.4}
\end{equation}
\begin{EESegment}{EE-TGT-P145-005}
Let \NSi{NS-F-P145-014}{\mathcal B_\rho} be the set of smooth coefficient sequences whose norm is
finite. This space is complete. In a Cauchy sequence, every coefficient and every order of its
derivative converges uniformly. Uniform convergence of one derivative after another then shows that
these limits are the derivatives of the limiting coefficient. In the rescaled variable
\NSi{NS-F-P145-015}{Y}, the logarithmic radial derivative is
\NSi{NS-F-P145-016}{\mathcal D_X=X\partial_X=Y\partial_Y}. Define
\NSi{NS-F-P145-017}{\mathcal I F=\int_0^Y F(Y',\eta)\,dY'}. Rescaling by
\NSi{NS-F-P145-018}{Y=\Lambda X} does not change radial averaging, so
\NSi{NS-F-P145-019}{\mathcal A_X(F)=Y^{-1}\mathcal I F}.
\end{EESegment}

\begin{lemma}\label{lem:B.1}
\begin{EESegment}{EE-TGT-P145-006}
Multiplication is bounded on \NSi{NS-F-P145-020}{\mathcal B_\rho}. For
\NSi{NS-F-P145-021}{\nu=1,2}, let \NSi{NS-F-P145-022}{\mathcal J_\nu F} be the solution
\NSi{NS-F-P145-023}{G} of
\NSi{NS-F-P145-024}{YG_{YY}+\nu G_Y=F} that extends smoothly to
\NSi{NS-F-P145-025}{Y=0} and satisfies \NSi{NS-F-P145-026}{G(0)=0}. Its coefficients are
\end{EESegment}
\NSFormula{NS-F-P145-027}{display}{B.5}
\begin{equation}
  (\mathcal J_\nu F)_{\alpha+1}
  =\frac{F_\alpha}{(\alpha+1)(\alpha+\nu)},
  \qquad (\mathcal J_\nu F)_0=0.
  \tag{B.5}\label{eq:B.5}
\end{equation}
\begin{EESegment}{EE-TGT-P145-007}
The operator \NSi{NS-F-P145-028}{\mathcal J_\nu} is bounded on
\NSi{NS-F-P145-029}{\mathcal B_\rho}. The same is true of radial averaging, multiplication by
\NSi{NS-F-P145-030}{Y}, and the operators \NSi{NS-F-P145-031}{\mathcal I} and
\NSi{NS-F-P145-032}{\partial_\eta\mathcal I}. Also,
\end{EESegment}
\NSFormula{NS-F-P145-033}{display}{B.6}
\begin{equation}
  \bigl\lVert\mathcal J_\nu[(\partial_\eta F)(\mathcal D_XG)]\bigr\rVert_\rho
  \leq C_\rho\lVert F\rVert_\rho\lVert G\rVert_\rho.
  \tag{B.6}\label{eq:B.6}
\end{equation}
\begin{EESegment}{EE-TGT-P145-008}
The same estimate still holds if either derivative is left out, or if \NSi{NS-F-P145-034}{F} is
replaced by \NSi{NS-F-P145-035}{\mathcal A_X(F)}. Extra factors without derivatives may also be
inserted, with one more algebra constant for each extra factor.
\end{EESegment}
\end{lemma}

\begin{proof}
\begin{EESegment}{EE-TGT-P145-009}
  The elementary convolution estimate
\NSFormula{NS-F-P145-036}{display}{B.7}
\begin{equation}
  \sum_{i=0}^{N}\frac{(N+1)^2}{(i+1)^2(N-i+1)^2}
  \leq 8\sum_{i=1}^{\infty}i^{-2}=:C_{\mathrm{sq}}
  \tag{B.7}\label{eq:B.7}
\end{equation}
  follows by splitting the sum at \NSi{NS-F-P145-037}{N/2}.
\end{EESegment}
\begin{EESegment}{EE-TGT-P145-010}
In the Leibniz formula for a product coefficient, the binomial coefficient from differentiation
cancels the factorials in the weights. The binomial ratio that remains is at most one, because
\end{EESegment}
\NSFormula{NS-F-P145-038}{display}{B.8}
\begin{equation}
  \binom{i+k}{k}\binom{\alpha-i+\beta-k}{\beta-k}
  \leq\binom{\alpha+\beta}{\beta}.
  \tag{B.8}\label{eq:B.8}
\end{equation}
\begin{EESegment}{EE-TGT-P145-011}
The reason is that the left side counts some, but not necessarily all, of the
\NSi{NS-F-P145-039}{\beta}-element subsets of a set split into blocks of sizes
\NSi{NS-F-P145-040}{i+k} and \NSi{NS-F-P145-041}{\alpha-i+\beta-k}. Apply
\eqref{eq:B.7} in both indices. This gives the algebra constant
\NSi{NS-F-P145-042}{C_{\mathrm{sq}}^2}.
\end{EESegment}

\NSPage{146}%
\begin{EESegment}{EE-TGT-P146-001}
The two weight ratios needed for integration are
\end{EESegment}
\NSFormula{NS-F-P146-001}{display}{B.9}
\begin{equation}
  \frac{\mathfrak a_{\alpha\beta}}{\mathfrak a_{\alpha+1,\beta}}
  =20\frac{(\alpha+2)^2}{(\alpha+1)^2}\frac{\alpha+1}{\alpha+\beta+1}
  \leq80,
  \tag{B.9}\label{eq:B.9}
\end{equation}
\NSFormula{NS-F-P146-002}{display}{B.10}
\begin{equation}
  \frac{\mathfrak a_{i,k+1}}{\mathfrak a_{i+1,k}}
  =\frac{20}{\rho}(i+1)\frac{(i+2)^2(k+1)^2}{(i+1)^2(k+2)^2}
  \leq\frac{80}{\rho}(i+1).
  \tag{B.10}\label{eq:B.10}
\end{equation}
\begin{EESegment}{EE-TGT-P146-002}
The first ratio proves the stated bounds for \NSi{NS-F-P146-003}{\mathcal J_\nu},
\NSi{NS-F-P146-004}{\mathcal I}, and multiplication by \NSi{NS-F-P146-005}{Y}. Averaging divides
the term of degree \NSi{NS-F-P146-006}{\alpha} by \NSi{NS-F-P146-007}{\alpha+1}. The second ratio
proves the bound for \NSi{NS-F-P146-008}{\partial_\eta\mathcal I}, because its integration divisor
\NSi{NS-F-P146-009}{i+1} cancels the matching factor.
\end{EESegment}

\begin{EESegment}{EE-TGT-P146-003}
For the mixed derivative bound \eqref{eq:B.6}, take the coefficient of degree
\NSi{NS-F-P146-010}{N=\alpha+1} on the left. Assign the extra degree to the factor with the parameter
derivative, and write \NSi{NS-F-P146-011}{I_1=i+1} and \NSi{NS-F-P146-012}{J_1=\alpha-i}.
After applying \eqref{eq:B.10} and dividing by \NSi{NS-F-P146-013}{\mathfrak a_{N\beta}}, the
\NSi{NS-F-P146-014}{\beta}th derivative is bounded by
\end{EESegment}
\NSFormula{NS-F-P146-015}{display}{none}
\begin{equation*}
  \frac{80}{\rho}\lVert F\rVert_\rho\lVert G\rVert_\rho
  \sum_{i=0}^{\alpha}\sum_{k=0}^{\beta}
  \frac{I_1J_1}{N(\alpha+\nu)}
  \frac{(N+1)^2}{(I_1+1)^2(J_1+1)^2}
  \frac{(\beta+1)^2}{(k+1)^2(\beta-k+1)^2}.
\end{equation*}
\begin{EESegment}{EE-TGT-P146-004}
Use \eqref{eq:B.8} with radial indices \NSi{NS-F-P146-016}{I_1,J_1} to drop the binomial ratio.
The first fraction is at most one, and \eqref{eq:B.7} bounds both sums. If
\NSi{NS-F-P146-017}{F} is averaged, its divisor \NSi{NS-F-P146-018}{i+1} cancels
\NSi{NS-F-P146-019}{I_1}, and the fraction that remains is again at most one. If the logarithmic
radial derivative is left out, replace \NSi{NS-F-P146-020}{J_1} by one. A logarithmic radial
derivative with no \NSi{NS-F-P146-021}{\eta}-derivative is controlled directly by
\eqref{eq:B.9} and \eqref{eq:B.5}. For extra factors, assign the indices in the same way, leave
those extra factors undifferentiated, and repeat the convolution estimate. This covers the endpoint
indices and every \NSi{NS-F-P146-022}{\beta}.
\end{EESegment}
\end{proof}

\begin{EESegment}{EE-TGT-P146-005}
\subsection{The nonlinear analytic profile at the axis.}\label{sec:B.3}
This subsection solves the equations \eqref{eq:4.13} for the axis data chosen in
Section~\ref{sec:B.1}, with the leading residual stress set to zero. Rescale the radius by
\NSi{NS-F-P146-023}{Y=\Lambda X}. The solution will be shown to approach an explicit profile as
\NSi{NS-F-P146-024}{\Lambda} grows, with the error stated in \eqref{eq:B.13}. The scalar comparison
function \NSi{NS-F-P146-025}{f_0} below controls the normalized azimuthal profile
\NSi{NS-F-P146-026}{\phi/\phi_*} through \NSi{NS-F-P146-027}{f_0(Y\chi)}. The bound
\eqref{eq:B.11} will show that it stays positive across the radial interval being used. That makes
division by \NSi{NS-F-P146-028}{\phi} valid in the profile equations.
\end{EESegment}

\begin{EESegment}{EE-TGT-P146-006}
The analytic scalar comparison function is
\end{EESegment}
\NSFormula{NS-F-P146-029}{display}{none}
\begin{equation*}
  f_0(z)=\sum_{\alpha\geq0}\frac{(-z/2)^\alpha}{\alpha!(\alpha+1)!}.
\end{equation*}
\begin{EESegment}{EE-TGT-P146-007}
For \NSi{NS-F-P146-030}{0\leq z\leq4.1}, set \NSi{NS-F-P146-031}{t=z/2}. After the cubic term,
the terms in the alternating tail get smaller in absolute value. It follows that
\end{EESegment}
\NSFormula{NS-F-P146-032}{display}{B.11}
\begin{equation}
  f_0(z)\geq1-\frac t2+\frac{t^2}{12}-\frac{t^3}{144}
  \geq\frac{305719}{1152000}>.265,
  \qquad f_0(z)\leq1.
  \tag{B.11}\label{eq:B.11}
\end{equation}
\begin{EESegment}{EE-TGT-P146-008}
The cubic decreases throughout \NSi{NS-F-P146-033}{[0,2.05]}, so its endpoint gives the value just
stated. For the upper bound, group the alternating tail after the constant term. Starting with the
linear term, its terms decrease in absolute value.
\end{EESegment}

\begin{proposition}\label{prop:B.2}
\begin{EESegment}{EE-TGT-P146-009}
Take the axis data from Section~\ref{sec:B.1}. There is a
\NSi{NS-F-P146-034}{\Lambda_0}. For every \NSi{NS-F-P146-035}{\Lambda\geq\Lambda_0}, there is an
amplitude threshold \NSi{NS-F-P146-036}{\mathsf C_0(\Lambda)} such that every
\NSi{NS-F-P146-037}{\mathsf C\geq\mathsf C_0(\Lambda)} gives an analytic axis profile whose leading
residual stress vanishes on \NSi{NS-F-P146-038}{0\leq Y=\Lambda X\leq4.1}. The profile has the form
\end{EESegment}
\NSFormula{NS-F-P146-039}{display}{B.12}
\begin{equation}
  \phi=\phi_*\Phi,
  \qquad U=U_*+\Lambda^{-1}\mathfrak u,
  \qquad \Pi=\Pi_0+\Lambda^{-1}\mathcal I(g^2\Phi^2),
  \qquad g=\phi_*/\mathsf C,
  \tag{B.12}\label{eq:B.12}
\end{equation}
\NSPage{147}%
\begin{EESegment}{EE-TGT-P147-001}
where \NSi{NS-F-P147-001}{\Phi(0,\eta)=1} and \NSi{NS-F-P147-002}{\mathfrak u(0,\eta)=0}. The
profile is analytic in \NSi{NS-F-P147-003}{Y} and \NSi{NS-F-P147-004}{\eta} on the same
neighborhood of \NSi{NS-F-P147-005}{[0,4.1]\times I}. For every fixed
\NSi{NS-F-P147-006}{r,s\geq0},
\end{EESegment}
\NSFormula{NS-F-P147-007}{display}{B.13}
\begin{equation}
  \sup_{[0,4.1]\times I}
  \left|\partial_Y^r\partial_\eta^s
    \left(\Phi-f_0(Y\chi),\mathfrak u+\frac{YZ_*}{2L}\right)\right|
  \leq\frac{C_{r,s}}{\Lambda}.
  \tag{B.13}\label{eq:B.13}
\end{equation}
\begin{EESegment}{EE-TGT-P147-002}
For all sufficiently large \NSi{NS-F-P147-008}{\Lambda} and all
\NSi{NS-F-P147-009}{\mathsf C\geq\mathsf C_0(\Lambda)}, the constants and a smaller common
parameter neighborhood do not depend on either large parameter. Written with derivatives in the
original radial variable, the bound is \NSi{NS-F-P147-010}{C_{r,s}\Lambda^{r-1}}.
\end{EESegment}
\end{proposition}

\begin{proof}
\begin{EESegment}{EE-TGT-P147-003}
Rewrite the equations using the rescaled radius. The nonlinear remainders then come with a factor
\NSi{NS-F-P147-011}{\Lambda^{-1}}. Applying the radial inverse operators turns the problem into a
contraction around the explicit comparison profile. In the notation of the stress equations, set
\end{EESegment}
\NSFormula{NS-F-P147-012}{display}{B.14}
\begin{equation}
  B=-2D_\eta\mathcal A_X(\mathfrak u)-d\partial_\eta\mathcal A_X(\mathfrak u),
  \qquad W=W_*+\Lambda^{-1}B,
  \qquad H_c=H_*+\Lambda^{-1}d\mathfrak u,
  \qquad \mathfrak p=\mathcal I(g^2\Phi^2).
  \tag{B.14}\label{eq:B.14}
\end{equation}
\begin{EESegment}{EE-TGT-P147-004}
The equations \eqref{eq:4.13} with zero leading residual stress, together with radial pressure
balance, are
\end{EESegment}
\NSFormula{NS-F-P147-013}{display}{B.15}
\begin{equation}
  -2L(X\phi_{XX}+2\phi_X)/\phi=S_q,
  \qquad -2L(XU_{XX}+U_X)=S_n,
  \qquad \Pi_X=\phi^2/\mathsf C^2.
  \tag{B.15}\label{eq:B.15}
\end{equation}
\begin{EESegment}{EE-TGT-P147-005}
Substitute \eqref{eq:B.12}. This gives
\end{EESegment}
\NSFormula{NS-F-P147-014}{display}{none}
\begin{align*}
  2(Y\Phi_{YY}+2\Phi_Y)&=-\chi\Phi+\Lambda^{-1}R_1,\\
  2(Y\mathfrak u_{YY}+\mathfrak u_Y)&=-Z_*/L+\Lambda^{-1}R_2,
\end{align*}
\begin{EESegment}{EE-TGT-P147-006}
where the remainders are exactly
\end{EESegment}
\NSFormula{NS-F-P147-015}{display}{none}
\begin{align*}
  LR_1={}&[W+h(1-2\eta U)+d\mathfrak u\zeta_*]\Phi+W\mathcal D_X\Phi+H_c\Phi_\eta,\\
  LR_2={}&[A(1-4\eta U_*)+dU_{*\eta}]\mathfrak u-2A_\eta\Lambda^{-1}\mathfrak u^2
            +W\mathcal D_X\mathfrak u+H_*\mathfrak u_\eta+d\Lambda^{-1}\mathfrak u\mathfrak u_\eta\\
          &\quad-4A_\eta\mathfrak p+d\partial_\eta\mathfrak p-2\eta\mathcal D_X\mathfrak p.
\end{align*}
\begin{EESegment}{EE-TGT-P147-007}
The leading angular term has this sign because
\NSi{NS-F-P147-016}{H_*\widetilde\xi_0/(\Lambda L)=-\chi}. To get the expanded axial remainder,
subtract the constant value \NSi{NS-F-P147-017}{-Z_*/L} from
\NSi{NS-F-P147-018}{L^{-1}\{A(1-2\eta U)U+W\mathcal D_XU+H_cU_\eta-4A_\eta\Pi+d\Pi_\eta-2\eta\mathcal D_X\Pi\}},
and multiply the result by \NSi{NS-F-P147-019}{\Lambda}.
\end{EESegment}

\begin{EESegment}{EE-TGT-P147-008}
The radial inverse bounds require control of the differentiated products in
\NSi{NS-F-P147-020}{R_1} and \NSi{NS-F-P147-021}{R_2}. The only products that contain both a
parameter derivative outside an integral and a logarithmic radial derivative are
\end{EESegment}
\NSFormula{NS-F-P147-022}{display}{none}
\begin{equation*}
  (\partial_\eta\mathcal A_X(\mathfrak u))\mathcal D_X\Phi,
  \qquad (\partial_\eta\mathcal A_X(\mathfrak u))\mathcal D_X\mathfrak u.
\end{equation*}
\begin{EESegment}{EE-TGT-P147-009}
The two derivatives fall on different factors. Also,
\NSi{NS-F-P147-023}{\partial_\eta\mathfrak p=\partial_\eta\mathcal I(g^2\Phi^2)} is bounded by
Lemma~\ref{lem:B.1}, and \NSi{NS-F-P147-024}{\mathcal D_X\mathfrak p=Yg^2\Phi^2}. It follows that
the maps \NSi{NS-F-P147-025}{(\Phi,\mathfrak u)\longmapsto\mathcal J_\nu R_i} are bounded and
locally Lipschitz on fixed balls in \NSi{NS-F-P147-026}{\mathcal B_\rho^2}. The displayed algebra
bounds give this quantitatively. To estimate the difference of two monomials, replace their factors
one at a time and subtract at each step.
\end{EESegment}

\begin{EESegment}{EE-TGT-P147-010}
These coefficient bounds are uniform in the large parameters. The functions
\NSi{NS-F-P147-027}{\zeta_*,\chi,U_*,Z_*/L} are fixed and holomorphic on
\NSi{NS-F-P147-028}{\Omega}. Choose a smaller neighborhood, and take
\NSi{NS-F-P147-029}{\rho} strictly smaller than its distance from the boundary of
\NSi{NS-F-P147-030}{\Omega}. After choosing \NSi{NS-F-P147-031}{\Lambda}, require
\end{EESegment}
\NSFormula{NS-F-P147-032}{display}{B.16}
\begin{equation}
  \mathsf C\geq\sup_{\eta\in\Omega}|\phi_*(\eta)|.
  \tag{B.16}\label{eq:B.16}
\end{equation}
\begin{EESegment}{EE-TGT-P147-011}
Then \NSi{NS-F-P147-033}{|g|\leq1} on \NSi{NS-F-P147-034}{\Omega}. Cauchy's inequality now
bounds its derivatives on the smaller neighborhood, uniformly in both
large parameters. Choosing \NSi{NS-F-P147-035}{\rho} strictly smaller absorbs the factor
\NSi{NS-F-P147-036}{(\beta+1)^2} in \eqref{eq:B.4}.
\end{EESegment}
\NSPage{148}%
\begin{EESegment}{EE-TGT-P148-001}
The bound on \NSi{NS-F-P148-001}{g} must hold on a complex neighborhood, because
\NSi{NS-F-P148-002}{\phi_*} itself can grow exponentially with
\NSi{NS-F-P148-003}{\Lambda}.
\end{EESegment}

\begin{EESegment}{EE-TGT-P148-002}
The order-one angular term still has to be inverted without assuming that its coefficient is small.
Let \NSi{NS-F-P148-004}{M_\chi} be a bound for multiplication by
\NSi{NS-F-P148-005}{\chi} on \NSi{NS-F-P148-006}{\mathcal B_\rho}, and set
\NSi{NS-F-P148-007}{T=\mathcal J_2\chi/2}, with multiplication done before integration. If
\NSi{NS-F-P148-008}{F} has no radial coefficient below degree \NSi{NS-F-P148-009}{b}, then
\eqref{eq:B.9} gives
\end{EESegment}
\NSFormula{NS-F-P148-010}{display}{none}
\begin{equation*}
  \lVert\mathcal J_2F\rVert_\rho
  \leq\frac{80}{(b+1)(b+2)}\lVert F\rVert_\rho,
  \qquad
  \lVert T^k\rVert\leq\frac{(40M_\chi)^k}{k!(k+1)!}.
\end{equation*}
\begin{EESegment}{EE-TGT-P148-003}
Multiplication by \NSi{NS-F-P148-011}{\chi} keeps every coefficient below degree
\NSi{NS-F-P148-012}{b} equal to zero, while \NSi{NS-F-P148-013}{\mathcal J_2} raises the lowest
possible nonzero degree by one. Parameter differentiation does not create any lower-degree
coefficients. The estimate therefore includes parameter derivatives of every order in these
iterated products. The series \NSi{NS-F-P148-014}{\sum_{k\geq0}(-T)^k} converges absolutely in
operator norm, and it is a two-sided inverse of \NSi{NS-F-P148-015}{1+T}.
\end{EESegment}

\begin{EESegment}{EE-TGT-P148-004}
The unperturbed solution is
\NSi{NS-F-P148-016}{\Phi^0=(1+T)^{-1}1=f_0(Y\chi)} and
\NSi{NS-F-P148-017}{\mathfrak u^0=-YZ_*/(2L)}. Solve the equations with the map
\end{EESegment}
\NSFormula{NS-F-P148-018}{display}{none}
\begin{equation*}
  (\Phi,\mathfrak u)\longmapsto
  \left(\Phi^0+\frac{(1+T)^{-1}\mathcal J_2R_1}{2\Lambda},
        \mathfrak u^0+\frac{\mathcal J_1R_2}{2\Lambda}\right).
\end{equation*}
\begin{EESegment}{EE-TGT-P148-005}
Take a fixed ball around \NSi{NS-F-P148-019}{(\Phi^0,\mathfrak u^0)}. For every sufficiently
large \NSi{NS-F-P148-020}{\Lambda}, the uniform bounds and Lipschitz constants above make the map
send this ball into itself, and make it a strict contraction there. Its fixed point is at most
\NSi{NS-F-P148-021}{C/\Lambda} away from the center in the norm. Choose
\NSi{NS-F-P148-022}{\mathsf C_0(\Lambda)} large enough to include \eqref{eq:B.16}. This keeps the
parameters in the order stated in the proposition.
\end{EESegment}

\begin{EESegment}{EE-TGT-P148-006}
For \NSi{NS-F-P148-023}{R<20},
\end{EESegment}
\NSFormula{NS-F-P148-024}{display}{none}
\begin{equation*}
  \sum_{\alpha\geq0}\binom{\alpha+\beta}{\beta}(R/20)^\alpha
  =(1-R/20)^{-\beta-1}.
\end{equation*}
\begin{EESegment}{EE-TGT-P148-007}
Choose \NSi{NS-F-P148-025}{4.1<R<20}. The coefficient norm then gives a single positive parameter radius that works
for all the functions involved. Any sufficiently small number below
\NSi{NS-F-P148-026}{\rho(1-R/20)} can be used for this radius on
\NSi{NS-F-P148-027}{|Y|\leq R}.
Cauchy estimates on a smaller \NSi{NS-F-P148-028}{Y} disk give \eqref{eq:B.13}, including every
fixed order of radial derivative. The derivatives in the original variable satisfy
\NSi{NS-F-P148-029}{\partial_X^r=\Lambda^r\partial_Y^r}. Because the data are real and the fixed
point is unique, the solution is real on the real rectangle. Analyticity in
\NSi{NS-F-P148-030}{X} makes the scalar profiles smooth at the axis. Along with
\NSi{NS-F-P148-031}{E=\sqrt{2X}\phi/\mathsf C}, Equations~\eqref{eq:4.4} to \eqref{eq:4.5}
then give a smooth Cartesian field. The lower bound \eqref{eq:B.11} gives
\NSi{NS-F-P148-032}{\Phi>0} when \NSi{NS-F-P148-033}{\Lambda} is large, and hence
\NSi{NS-F-P148-034}{\phi>0}. Division by \NSi{NS-F-P148-035}{\phi} in the first equation of
\eqref{eq:B.15} is therefore valid throughout the real rectangle.
\end{EESegment}
\end{proof}

\begin{EESegment}{EE-TGT-P148-008}
\subsection{Positive sources and the alternative at the outer endpoint.}\label{sec:B.4}
The analytic solution must satisfy the specific bounds needed to start its continuation toward
the outer profile. The approximation \eqref{eq:B.13} gives both a positive angular source
\NSi{NS-F-P148-036}{S_q} and the strict shear inequality \eqref{eq:B.19} at the end of the inner
interval. The same estimates will bound the normalized profiles and their parameter derivatives for
the joining construction in Proposition~\ref{prop:B.3} and Lemma~\ref{lem:B.4}.
\end{EESegment}

\begin{proposition}\label{prop:B.3}
\begin{EESegment}{EE-TGT-P148-009}
Increase \NSi{NS-F-P148-037}{\Lambda}, and then increase
\NSi{NS-F-P148-038}{\mathsf C_0(\Lambda)}, if needed. The normalized azimuthal profile satisfies
\NSi{NS-F-P148-039}{\Phi\geq c_0>0} on \NSi{NS-F-P148-040}{[0,4.1]\times[-1,1]}. On the same
rectangle, its angular source satisfies
\end{EESegment}
\NSFormula{NS-F-P148-041}{display}{B.17}
\begin{equation}
  S_q\geq2.5+.95\Lambda L\chi.
  \tag{B.17}\label{eq:B.17}
\end{equation}
\begin{EESegment}{EE-TGT-P148-010}
Let \NSi{NS-F-P148-042}{p_1=p_{s,1}=XQ_s/L},
\NSi{NS-F-P148-043}{p_2=p_{s,2}=XN_s/(LE)}, and \NSi{NS-F-P148-044}{n_s=N_s/L}. On
\NSi{NS-F-P148-045}{0<X\leq4.1/\Lambda},
\end{EESegment}
\NSFormula{NS-F-P148-046}{display}{B.18}
\begin{equation}
  p_1/X\geq c_1>0,
  \qquad 0<p_1\leq C_1,
  \qquad n_s=-2U_X=Z_*/L+O(\Lambda^{-1}).
  \tag{B.18}\label{eq:B.18}
\end{equation}
\NSPage{149}%
\begin{EESegment}{EE-TGT-P149-001}
At \NSi{NS-F-P149-001}{X_0=4/\Lambda}, there is a fixed
\NSi{NS-F-P149-002}{c_{\mathrm{ex}}>0} such that
\end{EESegment}
\NSFormula{NS-F-P149-003}{display}{B.19}
\begin{equation}
  p_1+\frac{p_2^2}{p_1}>2+c_{\mathrm{ex}}.
  \tag{B.19}\label{eq:B.19}
\end{equation}
\begin{EESegment}{EE-TGT-P149-002}
The quantities \NSi{NS-F-P149-004}{\Phi,\log\Phi,\mathfrak u,p_1}, and
\NSi{NS-F-P149-005}{n_s} have bounded derivatives at every fixed parameter order. Those bounds are
uniform in all later choices \NSi{NS-F-P149-006}{\mathsf C\geq\mathsf C_0(\Lambda)}. Constants for
derivatives in the original radial variable \NSi{NS-F-P149-007}{X} may depend on
\NSi{NS-F-P149-008}{\Lambda}.
\end{EESegment}
\end{proposition}

\begin{proof}
\begin{EESegment}{EE-TGT-P149-003}
The scalar bound \eqref{eq:B.11} and the approximation \eqref{eq:B.13} give
\NSi{NS-F-P149-009}{\Phi\geq c_0>0} when \NSi{NS-F-P149-010}{\Lambda} is large. On a smaller
common parameter neighborhood, \NSi{NS-F-P149-011}{\Phi} is also nonzero. This follows from
compactness and the uniform bound on its first derivative. It follows in turn that all parameter
derivatives of \NSi{NS-F-P149-012}{\log\Phi} are uniformly bounded.
\end{EESegment}

\begin{EESegment}{EE-TGT-P149-004}
The source bound uses the logarithmic derivatives of \NSi{NS-F-P149-013}{\Phi}. Since
\NSi{NS-F-P149-014}{|H_*\chi'|\leq2\lVert H_*'\rVert_\infty\chi} and
\NSi{NS-F-P149-015}{f_0} stays away from zero,
\NSi{NS-F-P149-016}{\mathcal D_X\log f_0(Y\chi)=O(\chi)} and
\NSi{NS-F-P149-017}{H_*\partial_\eta\log f_0(Y\chi)=O(\chi)}. The approximation and positivity
then give the following bound for the full solution.
\end{EESegment}
\NSFormula{NS-F-P149-018}{display}{B.20}
\begin{equation}
  |\mathcal D_X\log\Phi|+|H_*\partial_\eta\log\Phi|\leq C\chi+C/\Lambda.
  \tag{B.20}\label{eq:B.20}
\end{equation}
\begin{EESegment}{EE-TGT-P149-005}
Combine this with \eqref{eq:B.14} and \eqref{eq:B.3}. The result is
\end{EESegment}
\NSFormula{NS-F-P149-019}{display}{B.21}
\begin{align}
  -H_c\widetilde\xi_0&=\Lambda L\chi+O_{\sigma_*}(\sqrt\chi),\notag\\
  -H_c(\log\phi)_\eta
    &\geq\Lambda L\chi-C\chi-C_{\sigma_*}\sqrt\chi-C/\Lambda
      \geq.95\Lambda L\chi-C_{\sigma_*}/\Lambda.
  \tag{B.21}\label{eq:B.21}
\end{align}
\begin{EESegment}{EE-TGT-P149-006}
For the first error, use
\NSi{NS-F-P149-020}{|H_*|/(H_*^2+\sigma_*^2)\leq\sigma_*^{-1}\sqrt\chi}. For the second
inequality, increase \NSi{NS-F-P149-021}{\Lambda} until the term
\NSi{NS-F-P149-022}{C\chi} is absorbed. Then apply Young's inequality to the square-root term,
using a total of \NSi{NS-F-P149-023}{.05\Lambda L\chi}.
\end{EESegment}

\begin{EESegment}{EE-TGT-P149-007}
For the source itself, use \NSi{NS-F-P149-024}{l=1+\mathcal D_X\log\Phi} and the formula
\NSi{NS-F-P149-025}{S_q=-Wl-h(1-2\eta U)-H_c(\log E)_\eta}. The conditions
\NSi{NS-F-P149-026}{j_0\leq.05} and \eqref{eq:B.13} show, after increasing
\NSi{NS-F-P149-027}{\Lambda}, that \NSi{NS-F-P149-028}{|U|\leq4.1}. Hence
\NSi{NS-F-P149-029}{|h(1-2\eta U)|\leq10h}. The bounds above and
\NSi{NS-F-P149-030}{W-W_*=O(\Lambda^{-1})} now give
\end{EESegment}
\NSFormula{NS-F-P149-031}{display}{none}
\begin{equation*}
  S_q\geq\Lambda L\chi-W_*-10h-C\chi-C_{\sigma_*}\sqrt\chi-C/\Lambda.
\end{equation*}
\begin{EESegment}{EE-TGT-P149-008}
Because \NSi{NS-F-P149-032}{h\leq10^{-2}} and \NSi{NS-F-P149-033}{-W_*>2.8}, absorbing the terms in the same way
proves \eqref{eq:B.17} once \NSi{NS-F-P149-034}{\Lambda} is large. In particular, the
argument never divides by \NSi{NS-F-P149-035}{\chi} near a zero of
\NSi{NS-F-P149-036}{H_*}.
\end{EESegment}

\begin{EESegment}{EE-TGT-P149-009}
The function \NSi{NS-F-P149-037}{Q_s} extends smoothly to \NSi{NS-F-P149-038}{X=0}, and it has
the positive integral formula
\end{EESegment}
\NSFormula{NS-F-P149-039}{display}{none}
\begin{equation*}
  Q_s(X,\eta)
  =\frac{\displaystyle\int_0^X X'\Phi(\Lambda X',\eta)S_q(X',\eta)\,dX'}
         {X^2\Phi(\Lambda X,\eta)}.
\end{equation*}
\begin{EESegment}{EE-TGT-P149-010}
Uniform upper and lower bounds for \NSi{NS-F-P149-040}{\Phi}, together with \eqref{eq:B.17}, give
\NSi{NS-F-P149-041}{p_1/X\geq c_1}. Integrate the equations with zero leading residual stress from
\NSi{NS-F-P149-042}{X=0}, and use regularity at the axis. This gives
\end{EESegment}
\NSFormula{NS-F-P149-043}{display}{none}
\begin{equation*}
  p_1=a=-2Y\Phi_Y/\Phi,
  \qquad n_s=-2\mathfrak u_Y.
\end{equation*}
\begin{EESegment}{EE-TGT-P149-011}
The upper bound for \NSi{NS-F-P149-044}{p_1} and the approximation for
\NSi{NS-F-P149-045}{n_s} now follow from \eqref{eq:B.13}. The same identities give the stated
fixed bounds for the \NSi{NS-F-P149-046}{\eta}-derivatives.
\end{EESegment}

\begin{EESegment}{EE-TGT-P149-012}
It remains to prove the endpoint inequality. The two parameter regions chosen in \eqref{eq:B.2}
let either the azimuthal term or the axial term do the job. At \NSi{NS-F-P149-047}{Y=4}, first take
the case \NSi{NS-F-P149-048}{\chi>.99}. Then \NSi{NS-F-P149-049}{1.98<t=2\chi\leq2}. The
alternating series for \NSi{NS-F-P149-050}{f_0(z)+zf_0'(z)} gives
\end{EESegment}
\NSFormula{NS-F-P149-051}{display}{none}
\begin{equation*}
  f_0(z)+zf_0'(z)\leq1-t+\frac{t^2}{4}-\frac{t^3}{36}+\frac{t^4}{576}<-.18.
\end{equation*}
\begin{EESegment}{EE-TGT-P149-013}
This polynomial decreases on \NSi{NS-F-P149-052}{[1.98,2]}. At
\NSi{NS-F-P149-053}{t=1.98}, its value is
\NSi{NS-F-P149-054}{-75535511/400000000<-.188}. The decreasing alternating remainder has the
needed negative sign. By \eqref{eq:B.11},
\end{EESegment}
\NSPage{150}%
\begin{EESegment}{EE-TGT-P150-001}
\NSi{NS-F-P150-001}{-2zf_0'/f_0=2-2(f_0+zf_0')/f_0>2.36}. Equation \eqref{eq:B.13} then gives
\NSi{NS-F-P150-002}{p_1>2.3} at the outer endpoint when
\NSi{NS-F-P150-003}{\Lambda} is large.
\end{EESegment}

\begin{EESegment}{EE-TGT-P150-002}
In the other region, \NSi{NS-F-P150-004}{\chi\leq.99}, equation \eqref{eq:B.2} gives
\NSi{NS-F-P150-005}{|Z_*|>\delta_*}. Then \eqref{eq:B.18} gives
\NSi{NS-F-P150-006}{|n_s|\geq\delta_*/2} after increasing
\NSi{NS-F-P150-007}{\Lambda}. At the outer endpoint,
\end{EESegment}
\NSFormula{NS-F-P150-008}{display}{none}
\begin{equation*}
  |p_2|=\mathsf C\sqrt{2/\Lambda}\frac{|n_s|}{\phi_*\Phi}.
\end{equation*}
\begin{EESegment}{EE-TGT-P150-003}
Now \NSi{NS-F-P150-009}{\Lambda} is fixed. The denominator has a finite upper bound that is
uniform for all later large choices of \NSi{NS-F-P150-010}{\mathsf C}, while
\NSi{NS-F-P150-011}{p_1\leq C_1}. Increase \NSi{NS-F-P150-012}{\mathsf C_0(\Lambda)} until
\NSi{NS-F-P150-013}{p_2^2/p_1>2.3} throughout this region. Both cases give \eqref{eq:B.19}. One
possible choice is \NSi{NS-F-P150-014}{c_{\mathrm{ex}}=.2}.
\end{EESegment}
\end{proof}

\begin{EESegment}{EE-TGT-P150-004}
\subsection{A reference continuation with
\NSi{NS-F-P150-015}{p_{s,1}+p_{s,2}^2/p_{s,1}>2}.}\label{sec:B.5}
The next step continues the analytic axis profile from Proposition~\ref{prop:B.2} toward the outer
profile from Section~\ref{sec:A}. This continuation has to introduce a nonzero leading residual
stress \NSi{NS-F-P150-016}{\mathcal T_0} that satisfies the cone conditions, while keeping the
pressure datum \NSi{NS-F-P150-017}{\Pi_0} from Lemma~\ref{lem:A.5}. Start by constructing a
reference continuation \NSi{NS-F-P150-018}{\phi_{\mathrm r},U_{\mathrm r}}. The logarithmic radial
derivatives of \NSi{NS-F-P150-019}{\phi} and \NSi{NS-F-P150-020}{U} are gradually taken to zero,
as specified in \eqref{eq:B.22}. Lemma~\ref{lem:B.4} bounds the integrated residual coefficients
of this reference continuation. Those bounds will control the later change in shear in
Proposition~\ref{prop:B.5}.
\end{EESegment}

\begin{EESegment}{EE-TGT-P150-005}
As in Proposition~\ref{prop:B.3}, \NSi{NS-F-P150-021}{p_1,p_2} are the components of
\NSi{NS-F-P150-022}{p_s}, and \NSi{NS-F-P150-023}{n_s=N_s/L}. Recall that
\NSi{NS-F-P150-024}{\mathcal D_X=X\partial_X}. In a logarithmic radial coordinate, this is the same
as \NSi{NS-F-P150-025}{\partial_y}. A bound in \NSi{NS-F-P150-026}{C_\eta^k} concerns only
parameter derivatives, uniformly for \NSi{NS-F-P150-027}{-1\leq\eta\leq1}.
\end{EESegment}

\begin{EESegment}{EE-TGT-P150-006}
Set \NSi{NS-F-P150-028}{X_0=4/\Lambda}, \NSi{NS-F-P150-029}{X_{\mathrm b}=100}, and
\NSi{NS-F-P150-030}{X_i=110}. Choose a fixed \NSi{NS-F-P150-031}{\bar t>0} such that
\NSi{NS-F-P150-032}{4e^{2\bar t}<4.1}. Use the smooth step \NSi{NS-F-P150-033}{\sigma} from
\eqref{eq:A.5}. For \NSi{NS-F-P150-034}{0<t_1\leq\bar t}, set
\NSi{NS-F-P150-035}{y=\log(X/X_0)}. The reference profile agrees with the analytic axis profile
when \NSi{NS-F-P150-036}{y\leq t_1}. When \NSi{NS-F-P150-037}{t_1<y<2t_1}, prescribe
\end{EESegment}
\NSFormula{NS-F-P150-038}{display}{B.22}
\begin{align}
  \partial_y\log\phi_{\mathrm r}
    &=[1-\sigma((y-t_1)/t_1)]\mathcal D_X\log\phi_{\mathrm{nat}}(X,\eta),\notag\\
  \partial_yU_{\mathrm r}
    &=[1-\sigma((y-t_1)/t_1)]\mathcal D_XU_{\mathrm{nat}}(X,\eta),
  \tag{B.22}\label{eq:B.22}
\end{align}
\begin{EESegment}{EE-TGT-P150-007}
and then keep both fields constant in the logarithmic radial coordinate. Compute the reference
pressure and the functions \NSi{NS-F-P150-039}{Q_s,N_s} by integrating radially from the axis,
using the datum from Lemma~\ref{lem:A.5}.
\end{EESegment}

\begin{lemma}\label{lem:B.4}
\begin{EESegment}{EE-TGT-P150-008}
After fixing the parameters of the outer and axis profiles, first take
\NSi{NS-F-P150-040}{\Lambda} large enough, then take \NSi{NS-F-P150-041}{\mathsf C} large enough,
and finally take \NSi{NS-F-P150-042}{t_1} small enough. The following bounds then hold on
\NSi{NS-F-P150-043}{[X_0,X_i]\times[-1,1]}.
\end{EESegment}
\NSFormula{NS-F-P150-044}{display}{B.23}
\begin{equation}
  S_{q,\mathrm r}\geq.94L\Lambda\chi+2.4,
  \qquad l_{\mathrm r}\leq1,
  \qquad v_{\mathrm r}:=p_{1,\mathrm r}+\frac{p_{2,\mathrm r}^2}{p_{1,\mathrm r}}>2+c.
  \tag{B.23}\label{eq:B.23}
\end{equation}
\begin{EESegment}{EE-TGT-P150-009}
Here \NSi{NS-F-P150-045}{p_{1,\mathrm r}>0}, and
\NSi{NS-F-P150-046}{p_{1,\mathrm r}>3} whenever
\NSi{NS-F-P150-047}{X_{\mathrm b}\leq X\leq X_i}. For each fixed
\NSi{NS-F-P150-048}{k}, the parameter norms of \NSi{NS-F-P150-049}{\mathsf C E_{\mathrm r}}, its
reciprocal, \NSi{NS-F-P150-050}{p_{1,\mathrm r}}, and
\NSi{NS-F-P150-051}{n_{s,\mathrm r}} on \NSi{NS-F-P150-052}{[X_0,X_i]} are bounded independently
of all sufficiently large \NSi{NS-F-P150-053}{\mathsf C} and all
\NSi{NS-F-P150-054}{0<t_1\leq\bar t}.
\end{EESegment}
\end{lemma}

\begin{proof}
\begin{EESegment}{EE-TGT-P150-010}
Propositions~\ref{prop:B.2} and \ref{prop:B.3} show that the analytic axis profile has a positive
source and satisfies a strict inequality at its outer endpoint. Use the gradient bounds
\eqref{eq:B.20} and \eqref{eq:B.21}, together with the source bound \eqref{eq:B.17}, to carry these
properties through the cutoff. Integrating \eqref{eq:B.22} shows that
\NSi{NS-F-P150-055}{U_{\mathrm r}} and its radial average stay within
\NSi{NS-F-P150-056}{O(\Lambda^{-1})+o_{t_1}(1)} of \NSi{NS-F-P150-057}{U_*} in every required
parameter norm. This uses the fact that radial averaging is a contraction.
\end{EESegment}
\NSFormula{NS-F-P150-058}{display}{none}
\begin{equation*}
  \sup_{0<X\leq X_i}\lVert\mathcal A_X(U_{\mathrm r})-U_*\rVert_{C_\eta^k}
  \leq\sup_{0\leq X\leq X_i}\lVert U_{\mathrm r}-U_*\rVert_{C_\eta^k}.
\end{equation*}
\begin{EESegment}{EE-TGT-P150-011}
This estimate does not depend on the length of the interval where the reference profile is constant.
Integrating the same slopes also bounds
\NSi{NS-F-P150-059}{\log\phi_{\mathrm r}-\log\phi_*} in every parameter norm,
\end{EESegment}
